\documentclass[10pt,a4paper]{amsart}
\usepackage[margin=19mm]{geometry}
\usepackage{amsmath,amssymb,mathtools}
\usepackage[scr=rsfs,cal=euler]{mathalfa}
\usepackage{xfrac}
\usepackage[unicode,hidelinks,bookmarks=false]{hyperref}
\newcommand{\ie}{i.e.\ }
\newcommand{\cf}{cf.\ }
\newcommand{\resp}{respectively\ }
\newcommand{\Subsection}{\subsection}
\providecommand{\nobreakdash}{\nobreak\hbox{-}}
\setlength{\emergencystretch}{2em}
\multlinegap0pt
\usepackage{amsthm}
\usepackage{mathtools}
\usepackage{amsmath,amssymb,amsfonts}
\usepackage{graphicx}
\usepackage{bm}
\usepackage{enumitem}
\usepackage{mathrsfs}
\usepackage{xcolor}
\newtheorem{theorem}{Theorem}[section]
\newtheorem{prop}[theorem]{Proposition}
\newtheorem{lemma}[theorem]{Lemma}
\newtheorem{assumption}[theorem]{Assumption}
\newtheorem{corollary}[theorem]{Corollary}
\newtheorem{definition}{Definition}
\newtheorem{remark}{Remark}
\newtheoremstyle{named}{}{}{\bfshape}{}{\bfseries}#1{\thmnote{#3 }}
\theoremstyle{named}
\newtheorem*{namedtheorem}{Theorem}
\newtheoremstyle{named1}{}{}{\bfshape}{}{\bfseries}#1 {\thmnote{#3 }}
\theoremstyle{named1}
\newtheorem*{namedlemma}{Lemma}
\newtheorem{example}[theorem]{Example}
\theoremstyle{plain}
\newtheorem{thm}{\protect\theoremname}
\theoremstyle{remark}
\newtheorem{note}[thm]{\protect\notename}
\title[Classical sufficiency in quantum statistical experiments]{Classical Sufficiency in Quantum Statistical Experiments: the posterior covariance in the multinomial experiment (Lemma D.1)}
\author[S. Lahiry]{Samriddha Lahiry}
\date{Source: arXiv:2609.29348v1 (CC BY 4.0)}
\begin{document}
\maketitle
\vspace{1em}
\noindent\textit{Source and licence.} Classical Sufficiency in Quantum Statistical Experiments. Version arXiv:2609.29348v1 (CC BY 4.0), distributed under the Creative Commons Attribution 4.0 International licence, http://creativecommons.org/licenses/by/4.0/, as recorded in the arXiv OAI licence metadata for this exact version and shown on the abstract page https://arxiv.org/abs/2609.29348v1. This item is an editorial re-presentation of the paper's computation of the posterior covariance of the multinomial experiment under a Dirichlet prior on the ordered simplex, printed as Lemma D.1. A full rights notice is delivered at licenses/arxiv-2609.29348-CC-BY-4.0.txt.
\noindent\textit{Editorial scope.} Counted unit: the paper's computation of the posterior covariance of the multinomial experiment under a Dirichlet prior on the ordered simplex, printed as Lemma D.1 (source0.tex lines the statement at lines 6113--6168 and the proof at lines 6170--6635), delivered together with the paper's complete original proof of it as the single primary proof body. No mathematics is written, repaired or re-derived: statement and proof are the paper's own text line for line, in the paper's own notation. Every cross-reference inside the retained spans resolves inside the delivered document. Printed numbers are the paper's own: the wrapper restores the paper's counters before the retained material, so the counted statement prints as Lemma D.1, the paper's own number for it, because the wrapper restores the paper's appendix lettering and theorem counter before it. Omitted from this package and not redistributed: the paper's preamble and front matter as such (of which only the title and the author names are reproduced in the wrapper); the sections, lemmas, theorems and examples that lie outside the retained spans; and the paper's bibliography with its bib data file. No citation command occurs in any retained span, so the delivered document carries no bibliography and no reference or citation command of the delivered text was rewritten or adapted. Printed slips kept literal, among them the paper's own wording and its own choice of symbols. Layout and class: the delivered wrapper is the lane's pinned amsart editorial wrapper at 10pt on A4, to which this package adds the paper's own theorem-like declarations and the shorthand macros its retained text uses, all copied verbatim from the paper's source. No publisher class file, no publisher style file and no upstream script is redistributed. Assets: no retained span of this package contains a figure environment or a graphics-inclusion command, no image or artwork file exists in or is redistributed with this package, and the manifest and the delivery evidence both carry an empty asset-dependency list. A full rights notice is delivered at licenses/arxiv-2609.29348-CC-BY-4.0.txt.
\appendix
\setcounter{section}{4}
\setcounter{theorem}{0}
\begin{lemma}[Posterior covariance in the multinomial experiment]
\label{lem:multinomial-posterior-covariance}
Let $d\geq2$ be fixed, and define
\[
    \Theta_{\downarrow}
    :=
    \left\{
        \theta=(\theta_1,\ldots,\theta_d)^\top\in[0,1]^d:
        \theta_1\geq\cdots\geq\theta_d,\quad
        \sum_{i=1}^d\theta_i=1
    \right\}.
\]
Let
\[
    \Theta_{\mathrm{reg}}
    :=
    \left\{
        \theta\in\Theta_{\downarrow}:
        \theta_1>\cdots>\theta_d>0
    \right\}.
\]
Fix $\theta_0\in\Theta_{\mathrm{reg}}$. Write
\[
    \vartheta=(\theta_1,\ldots,\theta_{d-1})^\top,
    \qquad
    \theta_d=1-\sum_{i=1}^{d-1}\vartheta_i,
\]
and let $\vartheta_0$ denote the intrinsic coordinates of $\theta_0$.
Suppose that
\[
    N_n=(N_{n,1},\ldots,N_{n,d})
    \sim\operatorname{Mult}_d(n,\theta_0).
\]
Let $\Pi$ be a prior on $\Theta_{\downarrow}$ that is absolutely
continuous with respect to $(d-1)$-dimensional Lebesgue measure on the
affine hyperplane containing the simplex. Assume that its density has a
version that is continuous and strictly positive at $\theta_0$.
Define
\[
    V_{\vartheta_0}
    :=
    \operatorname{diag}(\theta_{0,1},\ldots,\theta_{0,d-1})
    -\vartheta_0\vartheta_0^\top.
\]
Then
\[
    \mathbb E_{\theta_0}
    \left[
        \left\|
            n\operatorname{Var}(\vartheta\mid N_n)
            -V_{\vartheta_0}
        \right\|_{\mathrm{op}}
    \right]
    \rightarrow0.
\]
\end{lemma}

\begin{proof}
Set $k:=d-1$ and define
\[
    \Omega_{\downarrow}
    :=
    \left\{
        \vartheta\in\mathbb R^k:
        \vartheta_1>\cdots>\vartheta_k>
        1-\sum_{i=1}^k\vartheta_i>0
    \right\}.
\]
The map
\[
    \theta(\vartheta)
    :=
    \left(
        \vartheta_1,\ldots,\vartheta_k,
        1-\sum_{i=1}^k\vartheta_i
    \right)^\top
\]
is an affine parametrization of $\Theta_{\mathrm{reg}}$.
Since $\Pi$ is absolutely continuous with respect to the Lebesgue measure, the complement $\Theta_{\downarrow}\setminus\Theta_{\mathrm{reg}}$
has $\Pi$-measure zero. Every posterior is absolutely
continuous with respect to $\Pi$; hence posterior integration may also be
restricted to $\Theta_{\mathrm{reg}}$.
Let $\widetilde\pi$ denote the induced prior density on
$\Omega_{\downarrow}$.

For $\vartheta\in\Omega_{\downarrow}$, the multinomial likelihood, up to
a factor independent of $\vartheta$, is
\[
    L_n(\vartheta)
    =
    \prod_{i=1}^k\vartheta_i^{N_{n,i}}
    \left(1-\sum_{i=1}^k\vartheta_i\right)^{N_{n,d}},
\]
with log-likelihood
\[
    \ell_n(\vartheta)
    =
    \sum_{i=1}^kN_{n,i}\log\vartheta_i
    +N_{n,d}\log\left(1-\sum_{i=1}^k\vartheta_i\right).
\]
Write
\[
    \widehat\theta_n:=N_n/n,
    \qquad
    \widehat\vartheta_n
    :=(N_{n,1}/n,\ldots,N_{n,k}/n)^\top.
\]

Since $\theta_0\in\Theta_{\mathrm{reg}}$,
\[
    \delta_0
    :=
    \min\left\{
        \theta_{0,d},
        \theta_{0,1}-\theta_{0,2},\ldots,
        \theta_{0,d-1}-\theta_{0,d}
    \right\}>0.
\]
Choose $r>0$ sufficiently small that
$\overline{B(\vartheta_0,3r)}\subset\Omega_{\downarrow}$ and
\[
    0<c_\pi\leq\widetilde\pi(\vartheta)\leq C_\pi<\infty,
    \qquad
    \vartheta\in\overline{B(\vartheta_0,3r)}.
\]
Set $U:=B(\vartheta_0,2r)$ and choose $\varepsilon_0>0$ such that
$4\varepsilon_0<\delta_0$ and $\sqrt{k}\,\varepsilon_0<r$.
Define
\[
    G_n
    :=
    \left\{
        \max_{1\leq i\leq d}
        |\widehat\theta_{n,i}-\theta_{0,i}|
        \leq\varepsilon_0
    \right\}.
\]
On $G_n$,
\[
    \widehat\theta_{n,d}
    \geq\theta_{0,d}-\varepsilon_0>3\varepsilon_0,
    \qquad
    \widehat\theta_{n,i}-\widehat\theta_{n,i+1}
    \geq\theta_{0,i}-\theta_{0,i+1}-2\varepsilon_0>0
\]
for $1\leq i<d$. Thus $\widehat\theta_n\in\Theta_{\mathrm{reg}}$ and
$\|\widehat\vartheta_n-\vartheta_0\|<r$ on $G_n$. In particular,
$\widehat\vartheta_n$ lies in a fixed compact subset of
$\Omega_{\downarrow}$. Since the empirical proportions maximize the
multinomial likelihood over the full simplex, they also maximize it over
$\Omega_{\downarrow}$ on this event. Hoeffding's inequality gives
\begin{equation}
    \mathbb P_{\theta_0}(G_n^c)
    \leq2d\exp(-2n\varepsilon_0^2).
    \label{eq:good-event-probability}
\end{equation}

All likelihood-ratio and information-matrix calculations at
$\widehat\vartheta_n$ below are performed on $G_n$. Auxiliary quantities
defined only on this event may be extended arbitrarily to $G_n^c$ when
stating convergence in probability, since
$\mathbb P_{\theta_0}(G_n^c)\to0$. The posterior distribution itself is
always the actual posterior, including on $G_n^c$.

For $\vartheta\in\Omega_{\downarrow}$, on $G_n$,
\[
\begin{aligned}
    \ell_n(\vartheta)-\ell_n(\widehat\vartheta_n)
    &=
    n\sum_{i=1}^d\widehat\theta_{n,i}
    \log\frac{\theta_i(\vartheta)}{\widehat\theta_{n,i}}\\
    &=-nD\!\left(
        \widehat\theta_n\,\middle\|\,\theta(\vartheta)
    \right),
\end{aligned}
\]
where $D(p\|q):=\sum_{i=1}^dp_i\log(p_i/q_i)$ uses natural logarithms.
Consequently,
\begin{equation}
    \frac{L_n(\vartheta)}{L_n(\widehat\vartheta_n)}
    =
    \exp\left\{
        -nD\!\left(
            \widehat\theta_n\,\middle\|\,\theta(\vartheta)
        \right)
    \right\}.
    \label{eq:likelihood-kl}
\end{equation}
This identity supplies both the local quadratic approximation and the
posterior tail bounds.

Introduce the local coordinate
$h:=\sqrt n(\vartheta-\widehat\vartheta_n)$ and its domain
\[
    \mathcal H_n
    :=
    \left\{
        h\in\mathbb R^k:
        \widehat\vartheta_n+h/\sqrt n\in\Omega_{\downarrow}
    \right\}.
\]
For every fixed $M<\infty$ and all sufficiently large $n$, on $G_n$,
\[
    \widehat\vartheta_n+h/\sqrt n\in U
    \quad\text{whenever }\|h\|\leq M.
\]
In particular, $\{h:\|h\|\leq M\}\subset\mathcal H_n$. Thus the
ordering constraints do not truncate the local limit around $\theta_0$.

Define the population Fisher information in intrinsic coordinates by
\[
    I(\vartheta)
    :=
    \operatorname{diag}\left(
        \frac1{\vartheta_1},\ldots,\frac1{\vartheta_k}
    \right)
    +\frac1{\theta_d(\vartheta)}\mathbf1_k\mathbf1_k^\top,
    \qquad
    \theta_d(\vartheta):=1-\sum_{i=1}^k\vartheta_i.
\]
The observed information at the empirical proportions satisfies
\begin{equation}
\begin{aligned}
    -\frac1n\nabla^2\ell_n(\widehat\vartheta_n)
    &=I(\widehat\vartheta_n)\\
    &=
    \operatorname{diag}\left(
        \frac1{\widehat\theta_{n,1}},\ldots,
        \frac1{\widehat\theta_{n,k}}
    \right)
    +\frac1{\widehat\theta_{n,d}}\mathbf1_k\mathbf1_k^\top.
\end{aligned}
    \label{eq:empirical-information}
\end{equation}
On the fixed neighborhood $B(\vartheta_0,3r)$, the third derivatives of
$n^{-1}\ell_n$ are uniformly bounded, independently of the sample.
Since $\nabla\ell_n(\widehat\vartheta_n)=0$ on $G_n$, Taylor's theorem
gives, for each fixed $M<\infty$ and all sufficiently large $n$,
\begin{equation}
    \sup_{\|h\|\leq M}
    \left|
        \ell_n(\widehat\vartheta_n+h/\sqrt n)
        -\ell_n(\widehat\vartheta_n)
        +\frac12h^\top I(\widehat\vartheta_n)h
    \right|
    \leq\frac{CM^3}{\sqrt n}
    \rightarrow0
    \label{eq:local-likelihood-expansion}
\end{equation}
on $G_n$. Moreover, $\widehat\vartheta_n\to\vartheta_0$ in
$\mathbb P_{\theta_0}$-probability, so
$I(\widehat\vartheta_n)\to I(\vartheta_0)$ in probability, where
\begin{equation}
    I(\vartheta_0)
    =
    \operatorname{diag}\left(
        \frac1{\theta_{0,1}},\ldots,\frac1{\theta_{0,k}}
    \right)
    +\frac1{\theta_{0,d}}\mathbf1_k\mathbf1_k^\top.
    \label{eq:fisher-information}
\end{equation}
Its inverse is
\begin{equation}
    I(\vartheta_0)^{-1}
    =V_{\vartheta_0}
    =\operatorname{diag}(\theta_{0,1},\ldots,\theta_{0,k})
     -\vartheta_0\vartheta_0^\top.
    \label{eq:inverse-information}
\end{equation}

For every sample, the posterior density on $\Omega_{\downarrow}$ is
\[
    \widetilde\pi_n(\vartheta\mid N_n)
    =
    \frac{L_n(\vartheta)\widetilde\pi(\vartheta)}{
        \displaystyle\int_{\Omega_{\downarrow}}
        L_n(u)\widetilde\pi(u)\,du}.
\]
Let $q_n(\cdot\mid N_n)$ denote the posterior density of
$H_n:=\sqrt n(\vartheta-\widehat\vartheta_n)$. On $G_n$, define
\[
    g_n(h)
    :=
    \exp\left\{
        \ell_n(\widehat\vartheta_n+h/\sqrt n)
        -\ell_n(\widehat\vartheta_n)
    \right\}
    \widetilde\pi(\widehat\vartheta_n+h/\sqrt n)
    1_{\mathcal H_n}(h),
\]
where the factors are evaluated only for $h\in\mathcal H_n$, and
$g_n(h)$ is defined to be zero otherwise. Then
\[
    q_n(h\mid N_n)=\frac{g_n(h)}{Z_n},
    \qquad
    Z_n:=\int_{\mathbb R^k}g_n(h)\,dh.
\]
The change of variables also gives, on $G_n$,
\begin{equation}
    \int_{\Omega_{\downarrow}}
    \frac{L_n(u)}{L_n(\widehat\vartheta_n)}
    \widetilde\pi(u)\,du
    =n^{-k/2}Z_n.
    \label{eq:intrinsic-posterior-normalizer}
\end{equation}
Write
\[
    \phi_{\vartheta_0}(h)
    :=
    \frac{\det\{I(\vartheta_0)\}^{1/2}}{(2\pi)^{k/2}}
    \exp\left\{-\frac12h^\top I(\vartheta_0)h\right\}
\]
for the density of $N_k(0,V_{\vartheta_0})$.

We first establish convergence on bounded sets. For every fixed
$M<\infty$, the likelihood expansion and convergence of the information
matrix yield
\begin{equation}
    \sup_{\|h\|\leq M}
    \left|
        e^{\ell_n(\widehat\vartheta_n+h/\sqrt n)
             -\ell_n(\widehat\vartheta_n)}
        -e^{-h^\top I(\vartheta_0)h/2}
    \right|
    \rightarrow0
    \label{local_expansion_exp}
\end{equation}
in $\mathbb P_{\theta_0}$-probability. Continuity of
$\widetilde\pi$ at $\vartheta_0$ also gives
\[
    \sup_{\|h\|\leq M}
    \left|
        \widetilde\pi(\widehat\vartheta_n+h/\sqrt n)
        -\widetilde\pi(\vartheta_0)
    \right|
    \rightarrow0
\]
in probability. Hence $g_n$ converges uniformly on every fixed bounded
set, in probability, to
\begin{equation}
    g(h):=\widetilde\pi(\vartheta_0)
    \exp\left\{-\frac12h^\top I(\vartheta_0)h\right\}.
    \label{eq:local-density-limit}
\end{equation}

We next control the posterior tails. Pinsker's inequality gives, for
$\vartheta\in\Omega_{\downarrow}$ and on $G_n$,
\[
    D\!\left(\widehat\theta_n\,\middle\|\,\theta(\vartheta)\right)
    \geq\frac12\|\widehat\theta_n-\theta(\vartheta)\|_1^2
    \geq\frac12\|\widehat\vartheta_n-\vartheta\|^2.
\]
Therefore
\begin{equation}
    \frac{L_n(\vartheta)}{L_n(\widehat\vartheta_n)}
    \leq
    \exp\left\{-\frac n2\|\vartheta-\widehat\vartheta_n\|^2\right\},
    \qquad \vartheta\in\Omega_{\downarrow},
    \label{eq:local-gaussian-bound}
\end{equation}
or, in local coordinates,
\begin{equation}
    \frac{L_n(\widehat\vartheta_n+h/\sqrt n)}{
          L_n(\widehat\vartheta_n)}
    \leq e^{-\|h\|^2/2},
    \qquad h\in\mathcal H_n.
    \label{eq:local-h-bound}
\end{equation}
Since $U=B(\vartheta_0,2r)$ and
$\|\widehat\vartheta_n-\vartheta_0\|<r$ on $G_n$,
\begin{equation}
    \inf_{\vartheta\in\Omega_{\downarrow}\setminus U}
    D\!\left(\widehat\theta_n\,\middle\|\,\theta(\vartheta)\right)
    \geq r^2/2=:c_U>0.
    \label{outside_U_bound}
\end{equation}
Consequently,
\begin{equation}
    \sup_{\vartheta\in\Omega_{\downarrow}\setminus U}
    \frac{L_n(\vartheta)}{L_n(\widehat\vartheta_n)}
    \leq e^{-nc_U}.
    \label{eq:outside-neighborhood-bound}
\end{equation}

A lower bound for $Z_n$ follows by restricting its integral to
$\|h\|\leq1$. For all sufficiently large $n$, on $G_n$, this ball is
mapped into $U$. The expansion \eqref{eq:local-likelihood-expansion},
the uniform upper bound for $I(\widehat\vartheta_n)$ on $G_n$, and the
local lower bound $c_\pi$ for the prior density therefore give a
deterministic constant $c_Z>0$ such that
\begin{equation}
    Z_n\geq c_Z
    \label{eq:normalizing-lower-bound}
\end{equation}
on $G_n$ for all sufficiently large $n$.

Split the posterior tail according to whether
$\widehat\vartheta_n+h/\sqrt n$ belongs to $U$. On the first part,
use $\widetilde\pi\leq C_\pi$ and \eqref{eq:local-h-bound}. On the
second part, use \eqref{eq:outside-neighborhood-bound}, boundedness of
the simplex, and
\[
    \int_{\mathcal H_n}
    \widetilde\pi(\widehat\vartheta_n+h/\sqrt n)\,dh
    =n^{k/2}.
\]
In particular, $1+\|h\|^2\leq Cn$ on $\mathcal H_n$. Together with
\eqref{eq:normalizing-lower-bound}, these estimates give, on $G_n$,
\begin{equation}
\begin{aligned}
    &\int_{\|h\|>M}(1+\|h\|^2)q_n(h\mid N_n)\,dh\\
    &\qquad\leq
    C\int_{\|h\|>M}(1+\|h\|^2)e^{-\|h\|^2/2}\,dh
    +Cn^{1+k/2}e^{-nc_U}
\end{aligned}
    \label{eq:posterior-tail-bound}
\end{equation}
for every $M>0$ and all sufficiently large $n$.
 The right-hand side
has the form $\varepsilon_M+r_n$, where $\varepsilon_M\downarrow0$
as $M\to\infty$ and $r_n\to0$ as $n\to\infty$. Thus the weighted
posterior tails vanish by first letting $n\to\infty$ and then
$M\to\infty$.

We now identify the normalization. Set
\[
    Z:=\int_{\mathbb R^k}g(h)\,dh
      =\widetilde\pi(\vartheta_0)(2\pi)^{k/2}
        \det\{I(\vartheta_0)\}^{-1/2}>0,
\]
so that $\phi_{\vartheta_0}=g/Z$. For fixed $M<\infty$, local uniform
convergence gives
\[
    A_{n,M}:=\int_{\|h\|\leq M}g_n(h)\,dh
    \rightarrow
    A_M:=\int_{\|h\|\leq M}g(h)\,dh
\]
in probability. By \eqref{eq:posterior-tail-bound}, on $G_n$,
\[
    \frac{A_{n,M}}{Z_n}
    =\int_{\|h\|\leq M}q_n(h\mid N_n)\,dh
    \geq1-\varepsilon_M-r_n.
\]
For $M$ sufficiently large that $\varepsilon_M<1/2$ and then $n$
sufficiently large, this yields
\[
    A_{n,M}\leq Z_n\leq\frac{A_{n,M}}{1-\varepsilon_M-r_n}.
\]
Since $A_M\uparrow Z$ and $\varepsilon_M\downarrow0$,
first letting $n\to\infty$ and then  $M\to\infty$, shows that
$Z_n\to Z$ in $\mathbb P_{\theta_0}$-probability.

It follows that, for every fixed $M$,
\[
    \int_{\|h\|\leq M}(1+\|h\|^2)
    |q_n(h\mid N_n)-\phi_{\vartheta_0}(h)|\,dh
    \rightarrow0
\]
in probability. On the complement,
\[
\begin{aligned}
    &\int_{\|h\|>M}(1+\|h\|^2)
      |q_n(h\mid N_n)-\phi_{\vartheta_0}(h)|\,dh\\
    &\qquad\leq
      \int_{\|h\|>M}(1+\|h\|^2)q_n(h\mid N_n)\,dh
      +\int_{\|h\|>M}(1+\|h\|^2)\phi_{\vartheta_0}(h)\,dh.
\end{aligned}
\]
The first term is at most $\varepsilon_M+r_n$ on $G_n$, and the second
tends to zero as $M\to\infty$. Since
$\mathbb P_{\theta_0}(G_n^c)\to0$, we conclude for the actual posterior
that
\begin{equation}
    \int_{\mathbb R^k}(1+\|h\|^2)
    |q_n(h\mid N_n)-\phi_{\vartheta_0}(h)|\,dh
    \rightarrow0
    \label{posterior-moment-conv}
\end{equation}
in $\mathbb P_{\theta_0}$-probability.

The weighted convergence \eqref{posterior-moment-conv} directly implies
convergence of the posterior first and second moments:
\[
    \mathbb E[H_n\mid N_n]\rightarrow0,
    \qquad
    \mathbb E[H_nH_n^\top\mid N_n]\rightarrow V_{\vartheta_0}
\]
in probability. Hence
\begin{equation}
    \operatorname{Var}(H_n\mid N_n)
    \rightarrow V_{\vartheta_0}
    \label{eq:posterior-covariance-probability}
\end{equation}
in probability.

To obtain convergence in expectation, the tail bound gives deterministic
constants $C<\infty$ and $n_0$ such that, for every $n\geq n_0$, on $G_n$,
\[
    \|\operatorname{Var}(H_n\mid N_n)\|_{\mathrm{op}}
    \leq\mathbb E[\|H_n\|^2\mid N_n]\leq C.
\]
On $G_n^c$, the posterior parameter $\vartheta$ lies in the bounded
ordered simplex and $\widehat\vartheta_n$ in the bounded full simplex.
Thus $\|H_n\|^2\leq Cn$ for every posterior draw, and
\[
    \mathbb E_{\theta_0}\left[
        \|\operatorname{Var}(H_n\mid N_n)\|_{\mathrm{op}}
        1_{G_n^c}
    \right]
    \leq Cn\,\mathbb P_{\theta_0}(G_n^c)
    \rightarrow0.
\]
Boundedness on $G_n$, convergence in probability, and the last estimate
therefore imply
\[
    \mathbb E_{\theta_0}\left[
        \|\operatorname{Var}(H_n\mid N_n)-V_{\vartheta_0}\|_{\mathrm{op}}
    \right]\rightarrow0.
\]
Finally, $\widehat\vartheta_n$ is fixed conditional on $N_n$, so
$\operatorname{Var}(H_n\mid N_n)=n\operatorname{Var}(\vartheta\mid N_n)$.
This proves the lemma.
\end{proof}

\end{document}
